\cotefile{75}
\scanpage{26}
c) Orientation d'un ensemble fini $I$\inserted{(cardinal $n$).} La donnée de $I$ équivaut à celle d'un torseur
\[
 P=\operatorname{Isom}(\deleted{\unsure{[1,n]_{\mathbb Z}}}I_n,I)=\operatorname{numérot}(I)
\]
$(I_n\mathrel{\unsure{:|=}}\{m\in\mathbb N^{*}\mid m\leq n\}=\{\deleted{\unsure{0}}1,2,\ldots,n\})$ sous le groupe
\[
 \deleted{\unsure{\mathfrak S|S}}\mathfrak S_n=\mathfrak S_{I_n}=\operatorname{Aut}(I_n).
\]

On a un hom canonique ``signature''
\diagram{\begin{tikzcd}[ampersand replacement=\&,column sep=large]
\mathfrak S_n \arrow[r,"s=\mathrm{sgn}"] \& \{\pm1\} \& \text{noyau }\mathfrak A_n
\end{tikzcd}}
et par définition l'ens des 2 orientations de $I$ est
\[
 \omega_I=\deleted{\unsure{\operatorname{Isom}}(}P\times_{\mathfrak S_n}\{\pm1\}\simeq P/\mathfrak A_n .
\]

\uline{Compatibilités avec} a) \uline{et} b)

\uline{Avec} b) On a sur $E_{(I,k)}=k^{I}$ un produit scalaire canonique
\[
 (x,y)=\sum_{i\in I}x_i y_i
\]
pour lequel $(e_i)$ est une base orthonormale. Ainsi $\deleted{E}I$ \deleted{\unsure{donne}}\inserted{\unsure{donne}} naissance à un module quadratique $E(I,\unsure{x|k|X})$. Ceci posé, on \inserted{définit} une bijection canonique
\[
 \omega_{E(I,x)}\simeq\omega_I
\]
\margin{Si $i_1,\mathinner{\cdotp\cdotp}i_n$ est une numérotation de $I$, $e_{i_1}\mathinner{\cdotp\cdotp}e_{i_n}$ est une base \inserted{\unsure{ord.}} orthon. de $E$, \ldots}
par transitivité dans
\diagram{\begin{tikzcd}[ampersand replacement=\&,column sep=large]
 \mathfrak S_n \arrow[r,hook,"\mathrm{incl.}"] \arrow[rr,bend right=24,"\mathrm{sgn}"] \& \mathrm{O}(\unsure{n|N},k) \arrow[r,"\det"] \& \{\pm1\}
\end{tikzcd}}

En particulier, on trouve (comp. avec a).),
\[
 \omega_{E(I,\mathbb R)\tikz[baseline=0pt]{\draw[->] (0,0) .. controls (.17,-.10) and (.31,-.22) .. (.49,-.22);}}\simeq\omega_I
\]
% The arrow starts at the closing parenthesis of E(I,R) and points down/right to the two-line note at the foot.
% The lower strokes of the second line meet the bottom edge of the scan.
\margin{considéré\\comme \unsure{vect.|variété} \unsure{réel|via|vie} \unsure{se pense|ses pts}}
% archive sheet number: 25
\scanpage{27}
\uline{Formules d'associativité}

a) Vectoriels réels. \unsure{Soit|Soient} $I$ ens fini, $(E_i)_{i\in I}$ famille de vectoriels \inserted{réels} de dim finie, $E=\coprod_{i\in I}E_i$ la somme directe ($\simeq$ produit direct). \deleted{Soit} $I^{+}$ $(I^{-})$ l'ens des $i\in I$ tels que $d_i=\operatorname{rg}_{\mathbb R}E_i$ soit pair (impair). On \unsure{a un} isom canonique
\[
 \boxed{\omega_E\simeq\omega_{I^{-}}\wedge\bigwedge_{i\in I}\omega_{E_i}}
\]
obtenu ainsi, en regardant le $2^{\text{ième}}$ mb comme. \deleted{\unsure{un}} quotient de
{\reversemarginpar\margin{\figure{\begin{tikzpicture}[scale=.34,transform shape]
\node[inner sep=2pt] (nb) {\parbox{5.9cm}{\uline{NB} Un ordre total de $I$ en induit un de $I^{-}$, donc
\diagram{\begin{tikzcd}[ampersand replacement=\&,column sep=small,row sep=-.3em]
\operatorname{numérot}(I)\arrow[r,"\text{épi}"] \& \phantom{\operatorname{numérot}(I^{-})} \\
\& \operatorname{numérot}(I^{-})
\end{tikzcd}}}};
\draw ([xshift=3pt]nb.north west) .. controls (nb.north) and ([xshift=-3pt]nb.north east) .. (nb.north east)
 .. controls ([xshift=-2pt]nb.east) and (nb.east) .. ([xshift=2pt]nb.south east)
 .. controls (nb.south east) and (nb.south) .. ([xshift=3pt]nb.south west);
\end{tikzpicture}}}\normalmarginpar}
\[
 \operatorname{numérot}(I^{-})\inserted{\times\operatorname{numérot}(I^{+})}\times\prod_{i\in I}\operatorname{bases}(E_i)
\]

Soit $\{i_1,\ldots,i_{n^{-}}\}$ num. de $I^{-}$, \inserted{$j_1\ \mathinner{\cdotp\cdotp}\ j_{n^{+}}$ une autre de $I^{+}$,} et $\forall\ i\in I$, $B_{\unsure{\alpha|i}}=(\deleted{B}e_{i,1},\mathinner{\cdotp\cdotp},e_{i,d_i})$ une base ordonnée de $E_i$, alors
\[
 e_{i_1,1}\unsure{2|,},\mathinner{\cdotp\cdotp},e_{i_1,d_1},e_{i_2,\unsure{2|1}},\mathinner{\cdotp\cdotp}\ e_{i_2,d_2},\mathinner{\cdotp\cdotp},e_{i_n,\unsure{t|1}},\mathinner{\cdotp\cdotp},e_{i_n,d_{i_n}}
\]
est une base de $E$. Si on change une $B_i$ par une matrice de $\mathrm{Gl}(d_i,\mathbb R)^{+}$, alors la base se trouve changée par une matrice de $\mathrm{Gl}(d,\mathbb R)^{+}$ $(d=\sum d_i\deleted{\unsure{)}}=\operatorname{rg}E)$, \inserted{donc définit $=$ él. de $\omega_E$.} Si on change l'ordre de $I^{-}$ par une permutation $\sigma^{-}$ (une transposition entre $\alpha$ et $\alpha+1$, disons \ldots) On change la base par une matrice de dét $\deleted{\unsure{(-1)^{i_{\alpha}i_{\alpha+1}}|(-1)^{i_{\alpha}+i_{\alpha+1}}}}\operatorname{sgn}\sigma^{-}$ ; par contre, si on change \deleted{par} l'ordre de $I^{+}$ par une $\sigma^{+}$, la matrice de passage est de det $+1$ en tous cas $\mathinner{\cdotp\cdotp\cdotp\cdotp}$
% archive sheet number: 26
\scanpage{28}
b) Contexte des familles finies de modules quadratiques : $=$ formule, avec $=$ \unsure{démonstration|démonstrati.}

c) \diagram{\begin{tikzcd}[ampersand replacement=\&]
\deleted{I}J \arrow[r,"p"] \& I
\end{tikzcd}}
application \inserted{(surjection \unsure{si on veut})} d'ens \deleted{\unsure{C|f}}finis,
\[
 I^{-}=\{i\in I\mid\operatorname{Card}p^{-1}(i)\ \text{impair}\}.
\]
\deleted{\unsure{Je dis que}} [p. ex $J$ ens avec relation d'équivalence \inserted{$\mathcal R$}, et $I=J/\mathcal R$]. \deleted{\unsure{O}}Alors
\[
 \omega_J\simeq\omega_{I^{-}}\wedge\bigwedge_{i\in I}\omega_{J_{\unsure{2|i}}}
\]
% slanted broken traces at the right; one estimated unreadable unit
\margin{\illegible{1}}

\uline{\unsure{Compatibilités|Compatibilité}} entre les isom \unsure{can.|cas.} a) b) c)

\figure{\begin{tikzpicture}
\draw (0,0) .. controls (.32,.03) and (.92,-.04) .. (1.2,-.01);
\end{tikzpicture}}

\uline{Cas d'extensions}
\[
 \{0\}=\deleted{\unsure{0}}E_0\subset E_1\subset E_2\ \mathinner{\cdotp\cdotp}\ \subset E_d=E
\]
\deleted{\unsure{On}} Définir un isom. canonique
\[
 \omega_E\simeq\bigwedge_{1\leq i\leq d}\omega_{E_i\unsure{/|{}}E_{i-1}}
\]
\figure{\begin{tikzpicture}
\draw (0,0) .. controls (.76,.06) and (1.95,.16) .. (2.90,.17);
\draw (1.04,-.05) .. controls (1.50,-.04) and (2.10,-.09) .. (2.46,-.05);
\draw[line cap=round] (-.01,-.07) -- (.04,-.085);
\end{tikzpicture}}
% archive sheet number: 27
\scanpage{30}
\figure{\begin{tikzpicture}
\node[anchor=north west,inner sep=0pt,outer sep=0pt] (letter) at (0,0) {\begin{minipage}{0.98\linewidth}
\footnotesize
% letterhead
\begin{minipage}[t]{0.56\linewidth}
\centering\sffamily
UNIVERSITE DES SCIENCES ET TECHNIQUES\\
DU LANGUEDOC\\
\rule{0.55cm}{0.4pt}\\
INSTITUT DE MATHEMATIQUES\\
Place Eugène Bataillon\\
\unsure{34060|34050} MONTPELLIER CEDEX\\
\rule{0.65cm}{0.4pt}\\
\end{minipage}\hfill
\begin{minipage}[t]{0.40\linewidth}
Montpellier, le\\[2.4\baselineskip]
% Two detached compact marks to the left of the typed dateline; origin unassigned.
\makebox[0pt][r]{\begin{tikzpicture}[x=1pt,y=1pt,baseline=0pt]
\fill (0,0.4) ellipse[x radius=0.35pt,y radius=0.6pt,rotate=-12];
\fill (1.45,1.4) ellipse[x radius=0.4pt,y radius=0.65pt,rotate=-12];
\end{tikzpicture}\hspace{1.6em}}Villecun le 30.01.78\\
\end{minipage}

\smallskip
Monsieur Alexandre Grothendieck\\
Professeur\\

\hfill\begin{minipage}{0.49\linewidth}
à Monsieur le Ministre de\\
\unsure{\deleted{l'Education Nationale}}\\
\end{minipage}

\smallskip
Monsieurle Ministre,\\

\deleted{Je me permets de vous envoyer copie d'une lettre datée de ce même jour, adressée à M. Rouzaud, Administrateur Provisoire de l'U.S.T.L., dans laquelle je lui soumets certains des problèmes de fonctionnement posés aux Universités par l'article 21 de l'ordonnance du 2.11.1945.}

\smallskip
\deleted{Par le même courrier j'envoie également copie de la même lettre à M. Peyrefitte, Ministre de la Justice, en attirant son attention sur les difficultés créées par cette ordonnance. Celles-ci ne touchent certes pas, à mon sens, l'aspect le plus important des problèmes soulevés par l'existence même ce cette ordonnance. Par contre, par leur nature technique, elles me semblent bien du ressort des responsabilités' officielles d'un Administrateur ou d'un Ministre. C'est pourquoi il m'a semblé opportu\deleted{\unsure{m|n}}\unsure{n|m} de vous \makebox[0pt][l]{\raisebox{0.65ex}{\tiny ç}}soumettre, en votre qualité de Ministre de l'Education Nationale,\mbox{\begin{tikzpicture}[x=1ex,y=1ex,baseline=0pt]
\fill (0.18,0.06) ellipse [x radius=0.07,y radius=0.055];
\fill (0.42,0.38) ellipse [x radius=0.07,y radius=0.055];
\end{tikzpicture}} des problèmes de votre compétence. Il me semblerait extrêmement souhaitable que vous preniez position publiquement à leur sujet.}

\smallskip
\hspace*{0.18\linewidth}\deleted{Veuillez agréer, Monsieur le Ministre, l'expression de}\\
\deleted{ma considération très distinguée .}\\

\hfill\begin{minipage}{0.31\linewidth}
\handwritten{A Grothendieck}\\[0.7\baselineskip]
% The paired detached strokes lie below and just to the right of the initial A.
\hspace*{0.7em}\begin{tikzpicture}[x=1pt,y=1pt,baseline=0pt]
\fill (0,0) .. controls (0.3,0.3) and (0.25,0.8) .. (0.45,1.05)
 .. controls (0.95,1.55) and (1.45,1.25) .. (1.45,0.9)
 .. controls (1.2,0.55) and (0.65,0.5) .. (0.45,0.45)
 .. controls (0.35,0.1) and (0.2,-0.15) .. (0,0);
\fill (3.6,-0.05) .. controls (3.9,0.1) and (3.85,0.55) .. (3.95,0.75)
 .. controls (4.25,0.95) and (4.7,1.0) .. (4.9,1.1)
 .. controls (5.2,1.15) and (5.25,0.75) .. (4.95,0.6)
 .. controls (4.6,0.5) and (4.25,0.4) .. (4.15,0.2)
 .. controls (4.1,-0.1) and (3.8,-0.3) .. (3.6,-0.05);
\end{tikzpicture}
\end{minipage}\\
\end{minipage}};
\begin{scope}[x={(letter.north east)},y={(letter.south west)}]
% Unassigned angular dark mark at the upper right physical edge; no letter or numeral supplied.
\draw[black,line width=0.55pt,line cap=round] (1.035,-0.026)
 .. controls (1.034,-0.022) and (1.034,-0.020) .. (1.035,-0.018)
 .. controls (1.037,-0.017) and (1.040,-0.017) .. (1.041,-0.018)
 -- (1.041,-0.020);
\fill (1.040,-0.023) ellipse[x radius=0.4pt,y radius=0.25pt];
\draw[gray,very thin] (0.654,0.258) -- (0.644,0.235)
 .. controls (0.561,0.272) and (0.327,0.389) .. (0.401,0.404)
 .. controls (0.449,0.420) and (0.738,0.382) .. (0.651,0.456)
 .. controls (0.579,0.511) and (0.399,0.568) .. (0.422,0.624)
 .. controls (0.451,0.668) and (0.719,0.641) .. (0.654,0.713)
 .. controls (0.593,0.781) and (0.477,0.842) .. (0.508,0.994);
\draw[gray,very thin] (0.654,0.254)
 .. controls (0.588,0.297) and (0.461,0.379) .. (0.443,0.407)
 .. controls (0.414,0.445) and (0.537,0.422) .. (0.528,0.449)
 .. controls (0.517,0.477) and (0.402,0.533) .. (0.492,0.583)
 .. controls (0.590,0.638) and (0.508,0.674) .. (0.475,0.733)
 .. controls (0.428,0.806) and (0.414,0.817) .. (0.431,0.889);
\end{scope}
\end{tikzpicture}}
% archive sheet number: 29
\scanpage{31}
% letterhead
\noindent\begin{minipage}[t]{0.57\linewidth}
\centering\rmfamily
UNIVERSITE DES SCIENCES ET TECHNIQUES\\
DU LANGUEDOC\\
\rule{0.55cm}{0.4pt}\\
\sffamily
INSTITUT DE MATHEMATIQUES\\
Place Eugène Bataillon\\
34060 MONTPELLIER CEDEX\\
\rule{0.65cm}{0.4pt}
\end{minipage}\hfill
\begin{minipage}[t]{0.37\linewidth}
Montpellier, le
\end{minipage}

\begin{center}
\setlength{\ULdepth}{1.45\baselineskip}
\uline{COURS\quad C4 1977-78 de $\mathbb A$. GROTHENDIECK}\par\vspace{.6\baselineskip}
\end{center}

\uline{Géométrie de l'icosaèdre}

Il est étrange que l'icosaèdre, étudié depuis plus de deux mille ans dans des contextes et des optiques très divers, ne semble jamais avoir été étudié systématiquement du point de vue combinatoire, ni de celui des réalisations géométriques sur des corps (voire des anneaux) généraux. Ce ``cours'' se voudrait une investigation commune sur ce thème, en partant de l'intuition combinatoire la plus naïve, exemplifiée sur des modèles en carton (avec un coloriage convenable des arètes).

Nous aborderons donc la théorie de l'i\deleted{e}cosaèdre d'un point de vue combinatoire d'abord, en explicitant dans le détail le dictionnaire entre la géométrie combinatoire de l'icosaèdre orienté, et celle de l'ensemble``orienté'' à cinq éléments (les cinq couleurs coloriant les arètes !) - ce qui donnera en même temps une motivation et une inspiration géométrique pour étudier le groupe des automorphismes commun aux deux situations, savoir le groupe alterné%
% Detached pair at approximately (1582,1659) and (1577,1666), between this source line and the next; hand and verbal role unassigned.
\margin{\figure{\begin{tikzpicture}[x=1mm,y=1mm]
\fill (0.60,0.90) .. controls (0.37,0.85) and (0.39,0.59) .. (0.57,0.58) .. controls (0.77,0.59) and (0.86,0.82) .. (0.60,0.90) -- cycle;
\fill (0.00,0.00) .. controls (-0.25,0.03) and (-0.24,-0.16) .. (-0.06,-0.22) .. controls (0.13,-0.21) and (0.22,-0.07) .. (0.00,0.00) -- cycle;
\end{tikzpicture}}}%
 $\handwritten{\unsure{\mathfrak A|\mathcal A}}_5$ (qui est aussi le plus petit groupe simple non abélien). La question de construire des ``réalisations géométriques'' de l'icosaèdre combinatoire, comme ``polyèdre régulier'' sur le corps des réels, ou sur tout autre corps, sera l'occasion à la fois d'approfondir la compréhension de la structure combinatoire, et de nous familiariser avec les notions de géométrie vectorielle et affine (voire de géométrie projective) sur des corps de caractéristique quelconque.

% Detached dark shape at the lower-left physical edge area, approximately [20,2293,49,2317]; hand and physical origin unassigned.
\par\vspace{2\baselineskip}
\noindent\hspace*{-1.5cm}\begin{minipage}{0.8cm}
\figure{\begin{tikzpicture}[x=1mm,y=1mm]
\fill (0.00,0.00) .. controls (-0.16,0.17) and (-0.03,0.40) .. (0.13,0.48)
 -- (0.17,0.72) -- (0.41,0.69) -- (0.50,1.05) -- (0.76,1.19)
 -- (0.98,1.08) -- (1.20,1.35) -- (1.48,1.48) -- (1.64,1.72)
 -- (1.91,1.86) -- (2.16,1.84) -- (2.42,2.02) -- (2.47,1.85)
 -- (2.71,1.74) -- (2.78,1.28) -- (2.50,1.25) -- (2.36,1.02)
 -- (2.08,1.09) -- (1.90,0.84) -- (1.70,0.82) -- (1.61,0.43)
 -- (1.37,0.32) -- (1.25,0.00) -- (0.95,-0.08) -- (0.73,-0.24)
 -- (0.50,-0.10) .. controls (0.24,-0.13) and (0.07,-0.12) .. (0.00,0.00) -- cycle;
\end{tikzpicture}}
\end{minipage}
% archive sheet number: 30
\scanpage{32}
% The regular F impressions and their lowered numerals are typed; the additional outer left strokes are handwritten.
% Unassigned compact hook above the text, source near (559,98).
\figure{\begin{tikzpicture}[x=0.006cm,y=-0.006cm,line cap=round,line join=round]
  \path[use as bounding box] (0,70) rectangle (1679,122);
  \draw[line width=0.3pt] (557.8,98.9) .. controls (560.3,99.3) and (560.8,98.0) .. (560.6,96.4);
\end{tikzpicture}}
En relation avec l'icosaèdre, nous découvrirons que les corps de caractéristique deux\unsure{.|} et cinq jouent un rôle très particulier, tant pour l'étude des réalisations géométriques, que pour l'étude combinatoire abstraite, qui peut s'expliciter par des notions de géométrie sur les corps ${\handwritten{\mathbb}{F}}_4$ et ${\handwritten{\mathbb}{F}}_5$ - ces deux dictionnaires étant subordonnés aux ``isomorphismes exceptionnels'' :
\[
 \begin{aligned}
 \handwritten{\unsure{\mathfrak A|\mathcal A}}_5 &\simeq \mathrm{Sl}\,(2,{\handwritten{\mathbb}{F}}_4)\quad(\simeq\mathrm{GP}(1,{\handwritten{\mathbb}{F}}_4)\simeq\mathrm{SO}(3,{\handwritten{\mathbb}{F}}_4))\\
 \handwritten{\mathfrak S}_5 &\simeq \mathrm{GP}\,(1,{\handwritten{\mathbb}{F}}_5)\quad(\simeq\mathrm{SO}(3,{\handwritten{\mathbb}{F}}_5))\ .
 \end{aligned}
\]
% Unassigned peripheral graphic traces. Their hand and textual role remain unassigned.
% Source-coordinate map; x increases rightward and y downward. Source dimensions: 1679 x 2362.
% Upper hook is drawn above the text. Lower map: right hook near (1469,1607); detached pair near (1504,1750);
% crooked left-edge trace near (41,2203); broken lower-central cluster near (1061,1985).
\figure{\begin{tikzpicture}[x=0.006cm,y=-0.006cm,line cap=round,line join=round]
  \path[use as bounding box] (0,1579) rectangle (1679,2233);
  \draw[line width=0.34pt] (1466.5,1609.0) .. controls (1469.9,1609.9) and (1471.3,1606.3) .. (1470.6,1604.5);
  \fill (1507.5,1746.3) ellipse[x radius=1.4,y radius=1.0];
  \fill (1502.5,1753.3) ellipse[x radius=1.5,y radius=0.9];
  \draw[black!50,line width=0.22pt] (38.7,2195.3) .. controls (38.8,2198.0) and (41.8,2199.0) .. (41.5,2201.5) .. controls (40.7,2205.1) and (42.8,2207.5) .. (44.1,2209.0);
  \fill[black!60] (44.1,2209.1) circle[radius=0.8];
  \draw[black!40,line width=0.18pt] (1059.4,1961.7) -- (1058.8,1964.8);
  \fill (1063.3,1974.2) circle[radius=0.75];
  \fill[black!55] (1057.0,1993.1) ellipse[x radius=1.1,y radius=0.7];
  \fill (1061.4,1995.6) ellipse[x radius=1.2,y radius=0.8];
  \draw[black!55,line width=0.18pt] (1064.0,2000.0) -- (1064.4,2002.4);
\end{tikzpicture}}
% archive sheet number: 31
\scanpage{33}
\begin{center}
Cours DEA 1977/78 de A. Grothendieck

\uline{Géométrie arithmétique des polygônes réguliers}
\end{center}

{\reversemarginpar\margin{\handwritten{!}}\normalmarginpar}%
% The hand of the marginal ! and short underline is unassigned.
% The short hand-drawn rule underlines the typed words en commun.
Ce ``cours'' se fera dans l'esprit d'une investigation \handwritten{\uline}{en commun} de multiples facettes de la notion de polygône régulier au-dessus d'un corps de base, puis d'un anneau de base quelconque. Ce sera une occasion stimulante pour nous familiariser avec les notions de base de géométrie vectorielle, \tikz[remember picture,overlay]{\coordinate (s033-affine);}affine, projective - avec ou sans métrique - en dimensions un, deux et trois, et avec les phénomènes particuliers à ces dimensions, ainsi qu'avec les aspects parfois imprévus que prennent des notions familières (comme celle \tikz[remember picture,overlay]{\coordinate (s033-de);}de réflexion\deleted{s}, ou de centre d'un polygône régulier) sur les corps de caractéristique positive, et notamment sur les corps finis. Le besoin d'une géométrie sur des anneaux de base généraux (tel l'anneau ``sur lequel s'écrit'' le $\mathrm{n}$-polygône régulier ``le plus général'') se fera sentir en cours de route, et nous amènera sans doute à affleurer quelques\\
\tikz[remember picture,overlay]{\coordinate (s033-points1);}points de fondements de géométrie algébrique. Si le temps le permet et si d'autres investigations ne nous retiennent , nous tenterons, dans la volée, de construire l'icosaèdre régulier ``général'' sur un anneau de base convenable, fini sur l'anneau \hspace{.5em}${\handwritten{\mathbb}{Z}}$\hspace{.5em} des entiers, ayant des corps résiduels de toutes caractéristiques. (La caractéristique deux \hspace{.5em}donnera sans doute du fil à retordre !).

A partir d'un point de départ intuitif et familier, ce ``cours'' voudrait\tikz[remember picture,overlay]{\coordinate (s033-voudrait);} être une ouverture à un univers fascinant, pleins de mystères qui ne demandent qu'à être élucidés. Il ne peut donner des fruits qu'avec une participation \tikz[remember picture,overlay]{\coordinate (s033-active);}active de l'étudiant et un travail personnel important, notamment pour développer a\deleted{e}vec\\
\tikz[remember picture,overlay]{\coordinate (s033-detail);}le détail nécessaire les nombreux points\handwritten{(}de fondement\inserted{\handwritten{s}} notamment) qui n'auront été qu'esquissés dans le ``cours''. Il lui est très fortement conseillé de suivre en même temps le cours \hspace{.5em}C4\hspace{.5em} sur l'icosaèdre, qui se fera dans un esprit analogue. Il y aura sans doute de nombreuses interférences entre ces deux cours de réflexion.

% The underlying e in avec has a visible mid-height horizontal bar.
% The live v is clear; the heavy overstrike leaves its inscription method open.
% The opening parenthesis retains C's hand encoding; reinforcement of type remains open.
% Z has two obliques; the added diagonal retains hand encoding at medium confidence,
% with a second type impression remaining an instrument alternative.
% C4 characters are clear; the tall narrow C and descending unfooted 4 leave
% inscription method unresolved. Plain encoding is retained without a glyph doubt.

% Detached geometries: no verbal role, origin or author is assigned.
% Local source-pixel geometry is scaled to the typeset x-height; the zero-width
% text anchors retain each mark's relation to the corresponding passage.
% The line break before le detail retains the left-edge mark's line attachment.
\figure{\begin{tikzpicture}[remember picture,overlay]
% Pale oval/U-shaped trace below the a of affine, full-source area x=1255--1264, y=638--650.
\begin{scope}[shift={(s033-affine)},x=0.26pt,y=-0.26pt]
\draw[black!30,line width=0.65pt,line cap=round] (3,0) .. controls (1,4) and (2,9) .. (6,10) .. controls (9,9) and (10,5) .. (8,2);
\end{scope}
% Compact point and paler broken trace below de before réflexion, source (1019,877).
\begin{scope}[shift={(s033-de)},x=0.26pt,y=-0.26pt]
\fill[black!45] (19,19) ellipse[x radius=1.25,y radius=.95];
\draw[black!25,line width=.45pt] (25.5,21.7) .. controls (27,20.9) and (29.5,21.3) .. (30.7,22.7);
\end{scope}
% Pale taper beside the source line beginning points, source (55--60,1220--1224).
\begin{scope}[shift={(s033-points1)},x=0.26pt,y=-0.26pt]
\fill[black!40] (-199,7) -- (-194.5,9.6) -- (-198.8,9.9) -- cycle;
\end{scope}
% Detached dark hook below the end of voudrait, around x=1498, y=1582.
\begin{scope}[shift={(s033-voudrait)},x=0.26pt,y=-0.26pt]
\draw[line width=0.75pt,line cap=round] (-17,16) .. controls (-15,20) and (-19,22) .. (-21,22);
\end{scope}
% Detached upper-right / lower-left pair below active, around (1534,1722), (1529,1729).
\begin{scope}[shift={(s033-active)},x=0.26pt,y=-0.26pt]
\fill (31,16) ellipse[x radius=1.5,y radius=1.1,rotate=45];
\fill (26,23) ellipse[x radius=1.3,y radius=1.0,rotate=20];
\end{scope}
% Compact two-lobed mark at the physical left edge beside le detail, around x=28--37,y=1839--1846.
\begin{scope}[shift={(s033-detail)},x=0.26pt,y=-0.26pt]
\fill (-227,-5) .. controls (-227,-8) and (-223,-8) .. (-222,-5)
  .. controls (-219,-6) and (-217,-5) .. (-217,-3)
  .. controls (-217,0) and (-220,1) .. (-221,0)
  .. controls (-223,0) and (-223,-3) .. (-224,-3)
  .. controls (-226,-3) and (-228,-3) .. (-227,-5) -- cycle;
\end{scope}
\end{tikzpicture}}
% archive sheet number: 32
